% ==========================================
% CAPÍTULO: CONCLUSIONES
% ==========================================

\chapter{Conclusiones}

% ==========================================
% CONCLUSIÓN - Bustillos Cruz Jonatan
% ==========================================
\newpage
\section{Bustillos Cruz Jonatan}

% TODO (Jonatan): Escribe tu conclusión sobre la práctica
% Incluye:
% - Qué aprendiste
% - Dificultades encontradas
% - Aplicaciones prácticas
% - Reflexión personal


\newpage
% ==========================================
% CONCLUSIÓN - Delgado Lucero Cristian Isaac
% ==========================================
\section{Delgado Lucero Cristian Isaac}

% TODO (Cristian): Escribe tu conclusión sobre la práctica
% Incluye:
% - Qué aprendiste
% - Dificultades encontradas
% - Aplicaciones prácticas
% - Reflexión personal



\newpage
% ==========================================
% CONCLUSIÓN - Frem Cortés José Angel
% ==========================================
\section{Frem Cortés José Angel}

% TODO (José Angel): Escribe tu conclusión sobre la práctica
% Incluye:
% - Qué aprendiste
% - Dificultades encontradas
% - Aplicaciones prácticas
% - Reflexión personal
A lo largo de esta práctica realicé un repaso general de los principales conceptos del lenguaje C, enfocándome tanto en los tipos de datos primitivos como en los derivados, lo cual me permitió reforzar conocimientos previamente adquiridos y comprender mejor su aplicación en programas reales. En particular, trabajé con tipos enteros, flotantes, caracteres y lógicos, observando sus características, diferencias de precisión, representación en memoria y formas de uso dentro de estructuras de control y operaciones básicas.

Uno de los aspectos más relevantes fue identificar cómo el lenguaje C maneja internamente ciertos tipos de datos, como el caso de los valores lógicos, que no cuentan con un tipo específico en su forma más básica, sino que se representan mediante enteros. Asimismo, analizar el uso de caracteres y su relación con la tabla ASCII me permitió entender mejor cómo se representan los datos a bajo nivel, al igual que reconocer las diferencias entre \texttt{float}, \texttt{double} y \texttt{long double}, lo cual me ayudó a comprender el impacto de la precisión en los cálculos.

En cuanto a los tipos derivados, desarrollé programas que involucraron arreglos, apuntadores, funciones y estructuras, lo que me permitió entender cómo organizar y manipular datos de manera más eficiente. Por ejemplo, el uso de arreglos me facilitó trabajar con múltiples datos del mismo tipo, mientras que los apuntadores me introdujeron al manejo directo de direcciones de memoria, siendo este uno de los temas que considero más complejos pero también más importantes dentro del lenguaje. Las funciones me permitieron estructurar mejor el código, dividiéndolo en partes más pequeñas y reutilizables, y las estructuras me ayudaron a agrupar diferentes tipos de datos en una sola entidad para representar información más completa.

Durante el desarrollo de los ejercicios, una de las principales dificultades fue mantener atención a los detalles sintácticos del lenguaje, como el uso correcto de operadores, especificadores de formato en \texttt{printf} y el manejo adecuado de apuntadores. Estos pequeños errores llegaron a provocar fallos en la ejecución o resultados incorrectos, por lo que me fue necesario ser cuidadoso y revisar constantemente el código. Sin embargo, estas dificultades también me permitieron aprender de los errores y mejorar mi lógica de programación.

Desde un punto de vista práctico, considero que los temas abordados tienen aplicaciones directas en el desarrollo de software, ya que constituyen la base para la construcción de programas más complejos. El correcto uso de tipos de datos, estructuras y funciones me permitió entender mejor cómo se desarrollan programas eficientes, especialmente en áreas relacionadas con el manejo de memoria y la organización de datos.

Además, esta práctica me permitió comprender el análisis y diseño de programas en paralelo utilizando procesos, lo cual amplía la forma en que se pueden desarrollar soluciones, ya que posibilita la ejecución simultánea de tareas. Esto me ayudó a entender la importancia de optimizar tiempos de ejecución y aprovechar mejor los recursos del sistema, así como su relación con el manejo de archivos y la interacción con el sistema operativo.

Como reflexión personal, este repaso me permitió consolidar mis conocimientos y detectar áreas en las que necesito seguir practicando, especialmente en el manejo de apuntadores y en la interpretación de errores. También comprendí que dominar los fundamentos es esencial para avanzar hacia temas más complejos. En general, considero que esta práctica fue útil para reforzar mi lógica de programación y aumentar mi confianza al trabajar con el lenguaje C.

\newpage
% ==========================================
% CONCLUSIÓN - Luna Gonzales Gabriel Alexis
% ==========================================
\section{Luna Gonzales Gabriel Alexis}

% TODO (Gabriel): Escribe tu conclusión sobre la práctica
% Incluye:
% - Qué aprendiste
% - Dificultades encontradas
% - Aplicaciones prácticas
% - Reflexión personal

